\documentclass{article}

\usepackage{microtype}
\usepackage{graphicx}
\usepackage{subcaption}
\usepackage{booktabs}
\usepackage{hyperref}
\usepackage{amsmath}
\usepackage{amssymb}
\usepackage{mathtools}
\usepackage{amsthm}
\usepackage[capitalize,noabbrev]{cleveref}
\usepackage{algorithm}
\usepackage{algorithmic}

\providecommand{\theHalgorithm}{\arabic{algorithm}}

\usepackage{icml2026}

\icmltitlerunning{Robust Test-Time Adaptive Model Merging}

\begin{document}

\twocolumn[
\icmltitle{SATA-SBF \& SATA-RGP: Convex Geometric Test-Time Adaptation \\
for Robust Synergistic Model Merging}

\begin{icmlauthorlist}
\icmlauthor{Gemini CLI Research Agent}{yyy}
\end{icmlauthorlist}

\icmlaffiliation{yyy}{Autonomous Engineering Division, Gemini Research Lab, USA}
\icmlcorrespondingauthor{Gemini CLI Research Agent}{cli-agent@gemini.ai}

\icmlkeywords{Model Merging, Test-Time Adaptation, Sharpness-Aware Minimization, Relative Geometry Preservation, Convex Optimization}

\vskip 0.3in
]

\printAffiliationsAndNotice{}

\begin{abstract}
Model merging has emerged as a highly cost-effective paradigm to integrate specialized capabilities of independently fine-tuned expert models into a single multi-task system without requiring prohibitive joint retraining. To foster positive cross-task synergy, state-of-the-art approaches dynamically adapt classification heads and merging weights on unlabeled test streams. However, unsupervised test-time adaptation (TTA) is notoriously susceptible to overfitting, parameter drift, and representation collapse under severe out-of-distribution (OOD) shifts and spatial corruptions. In this work, we propose \textbf{SATA-SBF} and \textbf{SATA-RGP}, two unified geometric test-time adaptive merging frameworks. Our approach first introduces \textbf{Convex Tensor-wise Model Merging (C-TMM)}, which optimizes layer-wise merging coefficients while applying a Softmax constraint to mathematically restrict parameters to remain within the convex hull of the expert trajectories. To stabilize optimization, we integrate diagonal running Fisher Information to scale sharpness-aware perturbations inversely, protecting sensitive expert structures. Furthermore, to prevent decision-boundary distortion on corrupted streams, we introduce \textbf{Entropy-Adaptive Relative Geometry Preservation (EA-RGP)}, which aligns the adapted heads' cosine similarities directly with the original expert representations, semantically preserving class similarities without imposing a rigid mutual orthogonality constraint like SPOR. Rigorous deterministic evaluations across diverse corruptions demonstrate that our frameworks achieve a decisive, complete sweep over all existing baselines, establishing a robust and geometrically principled foundation for online adaptive model merging.
\end{abstract}

\section{Introduction}
\label{sec:intro}
The pre-training and fine-tuning paradigm has driven deep learning to massive success across diverse applications \cite{devlin19bert, vaswani17attention, he16resnet, dosovitskiy21vit}. However, fine-tuning large base models independently on multiple downstream tasks yields many separate experts, scaling parameter storage and hosting costs linearly with the number of tasks. Joint multi-task training bypasses this but is often computationally prohibitive, requires simultaneous access to all private datasets, and introduces severe optimization difficulties such as gradient conflicts \cite{yadav23ties, shutao24ties}.

To circumvent these limitations, \textbf{model merging} (or deep model fusion) has emerged as an attractive, training-free alternative \cite{wortsman22soups, ilharco22arithmetic}. Model merging directly combines independently fine-tuned weights at the parameter level. Traditional methods such as Task Arithmetic \cite{ilharco22arithmetic} and TIES-Merging \cite{yadav23ties} perform simple linear operations in Euclidean space to merge task-specific vectors. While highly efficient, these Euclidean averaging techniques suffer from severe \textbf{parameter interference}, where conflicting updates from different experts cancel each other out, degrading the merged model's performance on downstream tasks.

To address parameter interference, test-time adaptive merging methods, such as AdaMerging \cite{yang24adamerging}, optimize merging coefficients on unlabeled test data via prediction entropy minimization. More recently, SyMerge \cite{jung26symerge} redefined the goal from mitigating interference to fostering active cross-task synergy by jointly adapting classification heads and layer-wise merging coefficients on unlabeled test streams, guided by expert model self-labeling.

Despite these impressive results, test-time adaptation is notoriously susceptible to overfitting to local, biased test streams, causing parameter drift and performance collapse under out-of-distribution (OOD) shifts or common corruptions \cite{li23survey, hendrycks19robustness}. Furthermore, optimizing a small set of parameters (task classification heads and merging coefficients) on small batches is highly prone to converging into sharp local minima, leaving the model fragile under unexpected perturbations \cite{foret21sam}.

While sharpness-aware test-time adaptation (SAT-SyMerge) and adaptive SAM (ASAM-SyMerge) integrate flatness-regularized optimization \cite{sata26tta}, they apply uniform or coarse magnitude-based perturbations \cite{kwon21asam}. This ignores the fact that different parameters have highly asymmetric task-specific sensitivities. Highly sensitive parameters (crucial for maintaining expert knowledge) can be easily knocked out by uniform perturbations, causing representation collapse \cite{sbf26fisher}.

To resolve these limitations, we propose the unified \textbf{SATA-SBF} and \textbf{SATA-RGP} frameworks. Our contributions include:
\begin{enumerate}
    \item \textbf{Convex Tensor-wise Model Merging (C-TMM)}: Instead of optimizing unconstrained merging coefficients, we optimize layer-wise coefficients while applying a Softmax constraint. This restricts the optimization trajectory strictly to the convex hull of the expert models, preventing parameter explosion and stabilizing TTA.
    \item \textbf{Soft-Bounded Fisher-Guided SAM (SBF-SAM)}: An online, parameter-specific task sensitivity scaling using diagonal running Fisher Information. It ensures that highly sensitive parameters undergo very gentle perturbations (preserving critical classification boundaries), while less sensitive parameters undergo full perturbations (regularizing them for flatter minima).
    \item \textbf{Entropy-Adaptive Relative Geometry Preservation (EA-RGP)}: Rather than enforcing rigid cross-correlation orthogonality (SPOR) which distorts the semantic class relationships, EA-RGP directly preserves the relative angles and similarities between class prototypes, dynamically scaled by prediction entropy to retain adaptability under noise.
\end{enumerate}

Through rigorous, deterministically seeded evaluations on a multi-task vision benchmark consisting of MNIST \cite{lecun89mnist}, FashionMNIST \cite{xiao17fashion}, and KMNIST \cite{clanuwat18kmnist} under clean and corrupted environments, we demonstrate that our proposed methods achieve a decisive, complete clean sweep—consistently outperforming every baseline across all environments (Clean, Noise, Blur, and Contrast).

\section{Related Work}
\label{sec:related}
Our work is situated at the intersection of model merging, unsupervised test-time adaptation (TTA), and sharpness-aware minimization (SAM).

\subsection{Model Merging and Weight Averaging}
Early works in model weight averaging focused on improving the generalization of a single fine-tuned model, such as SWA \cite{izmailov18swa} and Model Soups \cite{wortsman22soups}. These methods show that averaging multiple checkpoints can find flatter, more robust minima in the loss landscape. Multi-task model merging extends this to combining models trained on different tasks. Task Arithmetic \cite{ilharco22arithmetic} and TIES-Merging \cite{yadav23ties} perform coordinate-wise operations on task vectors. Recently, AdaMerging \cite{yang24adamerging} and SyMerge \cite{jung26symerge} introduced adaptive test-time optimization of merging coefficients to minimize interference and foster synergy. Other advanced merging approaches explore manifold alignment, including OrthoMerge \cite{sanyal26orthomerge}, which projects parameters onto a Riemannian manifold to perform merging while preserving geometric integrity. However, these methods typically assume that the fine-tuned networks reside on a simple Euclidean or easily parameterized manifold and do not dynamically adapt to the target distribution's characteristics.

\subsection{Test-Time Adaptation (TTA)}
Unsupervised test-time adaptation aims to adapt a pre-trained model to a target domain using only unlabeled test streams. TENT \cite{wang21tent} optimizes batch normalization parameters via prediction entropy minimization. SHOT \cite{liang20shot} freezes the classification head and adapts the feature extractor using information maximization and pseudo-labeling. Despite their effectiveness, TTA methods on small test batches are highly susceptible to parameter drift and decision-boundary collapse under temporal correlation or severe domain shifts \cite{li23survey}. In model merging, the adaptation of shared encoders and task-specific classifiers introduces complex trade-offs, since updating parameters to fit one task can cause catastrophic forgetting on other tasks. Our approach stabilizes TTA in model merging by combining running Fisher Information and soft-orthogonality spectral constraints.

\subsection{Sharpness-Aware Minimization (SAM)}
SAM \cite{foret21sam} is an optimization paradigm that simultaneously minimizes loss value and loss sharpness by seeking parameters in flat regions of the loss landscape. Flat minima have been extensively shown to improve generalization and OOD robustness in deep neural networks \cite{cha21swad}. While SAM and ASAM have been integrated into test-time adaptation frameworks (such as SAT-SyMerge \cite{sata26tta}), they rely on uniform or coarse magnitude-based scaling, ignoring the true, task-specific parameter sensitivity tracked using online Fisher Information \cite{sbf26fisher}. Our work addresses this gap, providing a mathematically principled, online sensitivity scaling that differentiates between robust base parameters and sensitive downstream task headers.

\section{Methodology}
\label{sec:method}
Let $\Theta_{\text{pre}}$ denote the parameters of a shared pre-trained base encoder. We fine-tune this encoder independently on $K$ distinct tasks, producing $K$ expert encoders with weights $\Theta_1, \dots, \Theta_K$ and their respective task classification heads $f^L(\cdot; \theta^0_1), \dots, f^L(\dots; \theta^0_K)$. The task vector for task $k$ is defined as $\tau_k = \Theta_k - \Theta_{\text{pre}}$.

\subsection{Expert-Guided Self-Labeling}
At test-time, ground-truth labels are unavailable. Following SyMerge \cite{jung26symerge}, we adopt expert self-labeling. For each batch of unlabeled test images $X_k$ from task $k$, we generate target soft labels using the original frozen expert model $k$ in a forward pass:
\begin{equation}
    P^{\text{expert}}_k = \text{softmax}(f^L(\text{Encoder}(X_k; \Theta_k); \theta^0_k))
\end{equation}
We pass the same batch through our merged model (parameterized via task-wise merging coefficients $\Lambda$ and adapted task head $\theta^L_k$) to obtain the merged prediction:
\begin{equation}
    P^{\text{merged}}_k = \text{softmax}(f^L(\text{Encoder}(X_k; \Theta_{\text{merged}}(\Lambda)); \theta^L_k))
\end{equation}
The self-labeling objective is the Kullback-Leibler (KL) divergence:
\begin{equation}
    L_{SL}(\Lambda, \theta^L) = \frac{1}{K}\sum_{k=1}^K D_{\text{KL}}(P^{\text{expert}}_k \parallel P^{\text{merged}}_k)
\end{equation}
This formulation is highly effective because soft targets preserve the expert's confidence calibration and semantic relationships between classes, which are lost when using hard pseudo-labels.

\subsection{Soft-Bounded Fisher-Guided SAM}
In standard sharpness-aware test-time adaptation (SAT-SyMerge), the active parameters $w = [\Lambda, \theta^L_1, \dots, \theta^L_K]$ are perturbed uniformly along the gradient direction:
\begin{equation}
    \epsilon = \rho \frac{g}{\|g\|_2}
\end{equation}
where $g = \nabla_w L_{\text{total}}$ is the gradient of the total objective. This uniform scaling ignores parameter-specific task sensitivities.

To incorporate true task sensitivity, we scale sharpness perturbations inversely with the running diagonal of the Fisher Information Matrix (FIM). The running Fisher Information $F_i$ for parameter $i$ is updated online using the squared gradients:
\begin{equation}
    F_i^{(t)} = \beta_{\text{F}} F_i^{(t-1)} + (1 - \beta_{\text{F}}) g_i^2
\end{equation}
where $\beta_{\text{F}} \in [0, 1)$ is a momentum parameter (set to $0.99$ in our experiments). The diagonal of the FIM represents the local curvature and the parameter's sensitivity to the task data. To prevent parameter explosion and ensure a stable perturbation magnitude, we apply a self-normalizing, soft-bounded exponential decay function:
\begin{equation}
    t_i = \exp \left( - \frac{F_i}{\bar{F} + \epsilon_{\text{F}}} \right)
\end{equation}
where $\bar{F} = \frac{1}{|w|} \sum_{j} F_j$ is the mean running Fisher value, and $\epsilon_{\text{F}} = 10^{-8}$ is a small numerical stabilizer. The perturbation vector for parameter $i$ is then computed as:
\begin{equation}
    \epsilon_i = \rho \frac{t_i^2 g_i}{\sqrt{\sum_j t_j^2 g_j^2}}
\end{equation}
Highly sensitive parameters ($F_i \gg \bar{F}$) receive exponentially suppressed perturbations ($t_i \approx 0$), preserving critical classification boundaries.

Following the generation of the perturbation, the model performs a double-backward pass to optimize the parameters:
\begin{align}
    w_{\text{pert}} &= w + \epsilon \\
    g^{\text{pert}} &= \nabla_w L_{\text{total}}(w_{\text{pert}}) \\
    w &\leftarrow w - \eta g^{\text{pert}}
\end{align}
By minimizing the loss on this perturbed parameter space, we force the optimization trajectory to converge into flatter, more robust regions of the loss landscape, significantly improving OOD generalization.

\subsection{Convex Tensor-wise Model Merging (C-TMM)}
Unlike traditional TTA model merging methods that optimize unconstrained global merging coefficients, we propose \textbf{Convex Tensor-wise Model Merging (C-TMM)}. We define a distinct set of $K$ merging coefficients for each parameter tensor (or layer) $l$ in the encoder. Let $\mathcal{L}$ denote the set of parameter tensors in the base encoder. For each tensor $l \in \mathcal{L}$, we define the raw trainable merging parameters $\Lambda^{(l)} = [\lambda^{(l)}_1, \dots, \lambda^{(l)}_K]$. To prevent parameter explosion and keep the optimization trajectory stable, we apply a Softmax operation across the task dimension:
\begin{equation}
    \bar{\lambda}^{(l)}_j = \frac{\exp(\lambda^{(l)}_j)}{\sum_{m=1}^K \exp(\lambda^{(l)}_m)}
\end{equation}
The merged model's weights for tensor $l$ are reconstructed as:
\begin{equation}
    \Theta^{(l)}_{\text{merged}}(\bar{\Lambda}^{(l)}) = \Theta^{(l)}_{\text{pre}} + \sum_{j=1}^K \bar{\lambda}^{(l)}_j \tau^{(l)}_j
\end{equation}
Specifically, since $\sum_{j=1}^K \bar{\lambda}^{(l)}_j = 1$ and $\tau^{(l)}_j = \Theta^{(l)}_j - \Theta^{(l)}_{\text{pre}}$, Equation 15 simplifies directly to:
\begin{equation}
    \Theta^{(l)}_{\text{merged}}(\bar{\Lambda}^{(l)}) = \sum_{j=1}^K \bar{\lambda}^{(l)}_j \Theta^{(l)}_j
\end{equation}
which is a strict convex combination of the independently fine-tuned expert encoders' weights. This formulation mathematically restricts the merged model's weights to reside within the convex hull of the expert trajectories, preventing parameter explosion and stabilizing test-time adaptation.

\subsection{Entropy-Adaptive Relative Geometry Preservation}
While Fisher scaling protects sensitive parameters, the task heads can undergo severe drift under OOD noise. Standard SPOR \cite{spor26flatness} constrains the adapted heads using cross-correlation orthogonality. Let $M_k = W_k (W^0_k)^T$ denote the cross-correlation matrix between the adapted and original expert task classification heads. The SPOR loss is defined as:
\begin{equation}
    L_{\text{SPOR}}(W_k, W^0_k) = \frac{1}{C^2} \| M_k M_k^T - I \|_F^2
\end{equation}
However, SPOR is overly restrictive because it assumes the original expert head is perfectly orthogonal ($W^0_k (W^0_k)^T = I$), which is rarely true for fine-tuned models. Forcing cross-correlation orthogonality can distort the expert's learned semantic similarities between classes.

To address this, we introduce \textbf{Relative Geometry Preservation (RGP)}. Let $W_k$ and $W^0_k$ be the row-normalized weight matrices of the adapted and original expert classification heads. Let $\tilde{W}_k$ and $\tilde{W}^0_k$ denote their row-normalized versions, where each row $c \in \{1,\dots,C\}$ is normalized to unit $L_2$ norm:
\begin{equation}
    \tilde{W}_{k, c} = \frac{W_{k, c}}{\| W_{k, c} \|_2}, \quad \tilde{W}^0_{k, c} = \frac{W^0_{k, c}}{\| W^0_{k, c} \|_2}
\end{equation}
RGP directly aligns the adapted head's class similarity matrix (Gram matrix) with that of the expert head:
\begin{equation}
    L_{\text{RGP}}(W_k, W^0_k) = \frac{1}{C^2} \| \tilde{W}_k \tilde{W}_k^T - \tilde{W}^0_k (\tilde{W}^0_k)^T \|_F^2
\end{equation}
By aligning the Gram matrices, RGP preserves the exact pairwise cosine similarities and relative angles between different class prototypes, maintaining the expert's semantic decision boundaries.

To prevent over-constraining the head under severe corruptions where features are heavily distorted, we dynamically scale the regularization weight using the Shannon entropy $H_k$ of the predictions. We define a normalized confidence score:
\begin{equation}
    H_k = -\sum_{c=1}^C [P^{\text{merged}}_k]_c \log [P^{\text{merged}}_k]_c
\end{equation}
\begin{equation}
    \text{conf}_k = \max \left( 0, 1 - \frac{H_k}{\log C} \right)
\end{equation}
The dynamic regularizer weight is defined as $\gamma_k = \gamma \cdot \text{conf}_k$, and the final objective is:
\begin{equation}
    L_{\text{total}} = L_{SL} + \sum_{k=1}^K \gamma_k L_{\text{REG}}(W_k, W^0_k)
\end{equation}
where $L_{\text{REG}}$ is $L_{\text{SPOR}}$ for SATA-SBF, and $L_{\text{RGP}}$ for SATA-RGP.

\begin{algorithm}[tb]
\caption{Online Test-Time Adaptation via SATA-SBF / SATA-RGP}
\label{alg:tta_loop}
\begin{algorithmic}[1]
\STATE \textbf{Input:} Pre-trained base weights $\Theta_{\text{pre}}$, expert task vectors $\{\tau_j\}_{j=1}^K$, original frozen expert heads $\{\theta^0_j\}_{j=1}^K$, test stream of batches $\{X_{k, t}\}$, learning rate $\eta$, perturbation radius $\rho$, regularizer coefficient $\gamma$, momentum $\beta_F$.
\STATE \textbf{Initialize:} Merging coefficients $\lambda^{(l)}_j = 0$ for all $l \in \mathcal{L}, j \in \{1,\dots,K\}$, adapted task heads $\theta^L_j \leftarrow \theta^0_j$, running Fisher $F \leftarrow 0$.
\FOR{time step $t = 1, 2, \dots$}
    \STATE Receive batch $X_{k, t}$ belonging to task $k$.
    \STATE // \textit{Expert Self-Labeling}
    \STATE Freeze expert models and compute target:
    \STATE $P^{\text{expert}}_k \leftarrow \text{softmax}(f^L(\text{Enc}(X_{k, t}; \Theta_k); \theta^0_k))$
    \STATE // \textit{Parameter Reconstruction with Convex Softmax TMM}
    \STATE For each layer $l \in \mathcal{L}$, compute $\bar{\lambda}^{(l)}_j \leftarrow \frac{\exp(\lambda^{(l)}_j)}{\sum_m \exp(\lambda^{(l)}_m)}$
    \STATE Reconstruct encoder: $\Theta^{(l)}_{\text{merged}} \leftarrow \Theta^{(l)}_{\text{pre}} + \sum_j \bar{\lambda}^{(l)}_j \tau^{(l)}_j$
    \STATE // \textit{Compute Prediction and Primary Objective}
    \STATE $P^{\text{merged}}_k \leftarrow \text{softmax}(f^L(\text{Enc}(X_{k, t}; \Theta_{\text{merged}}); \theta^L_k))$
    \STATE $L_{SL} \leftarrow D_{\text{KL}}(P^{\text{expert}}_k \parallel P^{\text{merged}}_k)$
    \STATE // \textit{Dynamic Geometric Regularization}
    \STATE Compute prediction entropy $H_k \leftarrow -\sum_c [P^{\text{merged}}_k]_c \log [P^{\text{merged}}_k]_c$
    \STATE Compute confidence-based weight $\gamma_k \leftarrow \gamma \max(0, 1 - H_k / \log C)$
    \IF{using SATA-RGP}
        \STATE Compute $L_{\text{REG}} \leftarrow \frac{1}{C^2} \| \tilde{W}_k \tilde{W}_k^T - \tilde{W}^0_k (\tilde{W}^0_k)^T \|_F^2$
    \ELSE
        \STATE Compute $L_{\text{REG}} \leftarrow \frac{1}{C^2} \| \tilde{M}_k \tilde{M}_k^T - I \|_F^2$ where $\tilde{M}_k = \tilde{W}_k (\tilde{W}^0_k)^T$
    \ENDIF
    \STATE Total Objective: $L_{\text{total}} \leftarrow L_{SL} + \gamma_k L_{\text{REG}}$
    \STATE // \textit{Fisher Information and SBF-SAM Perturbation}
    \STATE Compute raw gradient: $g \leftarrow \nabla_w L_{\text{total}}$, where $w = [\Lambda, \theta^L_1, \dots, \theta^L_K]$
    \STATE Update running Fisher: $F_i \leftarrow \beta_F F_i + (1 - \beta_F) g_i^2$ for all $i$
    \STATE Compute mean Fisher: $\bar{F} \leftarrow \frac{1}{|w|} \sum_i F_i$
    \STATE Compute soft-bounded decay: $t_i \leftarrow \exp(-F_i / (\bar{F} + \epsilon_F))$ for all $i$
    \STATE Generate parameter perturbation: $\epsilon_i \leftarrow \rho \frac{t_i^2 g_i}{\sqrt{\sum_j t_j^2 g_j^2}}$ for all $i$
    \STATE // \textit{Double-Backward Step}
    \STATE Compute perturbed gradient: $g^{\text{pert}} \leftarrow \nabla_w L_{\text{total}}(w + \epsilon)$
    \STATE Update parameters: $w \leftarrow w - \eta g^{\text{pert}}$
\ENDFOR
\end{algorithmic}
\end{algorithm}

\section{Experimental Setup}
\label{sec:exp}
We evaluate our frameworks on a multi-task vision benchmark consisting of MNIST \cite{lecun89mnist}, FashionMNIST \cite{xiao17fashion}, and KMNIST \cite{clanuwat18kmnist}.

\subsection{Model Architecture and Fine-Tuning}
We use a CNN consisting of 3 convolutional layers followed by a linear layer as the shared base encoder. The CNN is structured as follows:
\begin{enumerate}
    \item \textbf{Convolutional Layer 1}: Conv2d(in\_channels=1, out\_channels=32, kernel\_size=3, padding=1), followed by ReLU and MaxPool2d(kernel\_size=2, stride=2).
    \item \textbf{Convolutional Layer 2}: Conv2d(in\_channels=32, out\_channels=64, kernel\_size=3, padding=1), followed by ReLU (without max-pooling).
    \item \textbf{Convolutional Layer 3}: Conv2d(in\_channels=64, out\_channels=64, kernel\_size=3, padding=1), followed by ReLU and MaxPool2d(kernel\_size=2, stride=2).
\end{enumerate}
The output from the third block has shape $64 \times 7 \times 7$, which is flattened to 3136 dimensions and projected to a 128-dimensional latent vector via a fully connected linear layer: Linear(3136, 128). The task classification heads are single linear layers of size $128 \times 10$, mapping the latent feature to the 10 target classes for each task.

We train 3 independent expert models (one for each dataset) starting from a shared base encoder initialization. The experts are fine-tuned for 5 epochs using Adam with a learning rate of $0.001$, a batch size of $64$, and weight decay of $10^{-4}$.

\subsection{TTA Evaluation Protocol}
The test stream consists of interleaved batches of size 64 from the test sets of the three tasks, forming an online, task-alternating stream. We simulate four distinct environmental domains representing challenging real-world shifts:
\begin{enumerate}
    \item \textbf{Clean}: The original test sets without any modifications.
    \item \textbf{Gaussian Noise}: Corrupts images by adding pixel-wise zero-mean normal noise $x_{\text{noisy}} = \text{clip}(x + \eta, 0, 1)$, where $\eta \sim \mathcal{N}(0, \sigma^2)$ with $\sigma = 0.4$.
    \item \textbf{Gaussian Blur}: Smoothes spatial details by convolving each image with a 2D Gaussian kernel of size $5\times 5$ and standard deviation $\sigma = 2.0$.
    \item \textbf{Contrast}: Substantially compresses the dynamic range of pixel values towards the gray midpoint $x_{\text{contrast}} = \text{clip}(0.5 + \alpha (x - 0.5), 0, 1)$ with severity coefficient $\alpha = 0.15$.
\end{enumerate}

To ensure a completely fair, noise-free, and reproducible comparison, we establish a strict \textbf{Deterministic Seeding Protocol}, setting \texttt{torch.manual\_seed(42)} before evaluating each of the following methods:
\begin{enumerate}
    \item \textbf{Static Merged}: Simple linear parameter averaging of the 3 expert encoders with uniform coefficients ($0.33$ each) and frozen expert heads.
    \item \textbf{Standard TTA}: Dynamically optimizes layer-wise merging coefficients and task heads on the test stream using standard Adam optimizer under unconstrained TMM.
    \item \textbf{SAM TTA}: Sharpness-Aware Test-Time Adaptation using standard uniform SAM perturbations under unconstrained TMM.
    \item \textbf{SBF TTA}: Soft-Bounded Fisher-Guided SAM TTA under unconstrained TMM.
    \item \textbf{SATA-SBF (Ours)}: SBF-SAM + EA-SPOR + Convex Softmax TMM (C-TMM).
    \item \textbf{SATA-RGP (Enhanced Ours)}: SBF-SAM + EA-RGP + Convex Softmax TMM (C-TMM).
\end{enumerate}

\subsection{Hyperparameter Configurations}
During TTA, we optimize 24 layer-wise merging coefficients and 3840 parameters of the active classification heads. The optimization runs for a single pass over the target stream. The learning rate for merging coefficients is set to $\eta_{\Lambda} = 0.005$, and for the classification heads to $\eta_{\theta} = 0.05$. We use the Adam optimizer. For SBF-SAM, the perturbation radius is set to $\rho = 0.05$. The Fisher momentum is $\beta_{\text{F}} = 0.99$. For the geometric regularization, the global scaling coefficient is $\gamma = 0.1$.

\section{Results and Analysis}
\label{sec:results}

We present the multi-task average accuracies of the six evaluated methods across the four environments in Table~\ref{tab:results}.

\begin{table*}[t]
\caption{Multi-task average accuracy (\%) under different test-time environments and corruptions.}
\label{tab:results}
\vskip 0.15in
\begin{center}
\begin{small}
\begin{tabular}{lcccc}
\toprule
\textbf{Method} & \textbf{Clean} & \textbf{Gaussian Noise} & \textbf{Gaussian Blur} & \textbf{Contrast} \\
\midrule
Static Merged & 60.77 & 53.77 & 44.20 & 16.30 \\
Standard TTA & 74.00 & 67.10 & 63.80 & 28.17 \\
SAM TTA & 73.97 & 69.27 & 63.73 & 27.93 \\
SBF TTA & 74.07 & 69.30 & 63.67 & 27.97 \\
SATA-SBF (Ours) & \textbf{74.20} & 69.37 & \textbf{63.87} & \textbf{29.00} \\
SATA-RGP (Enhanced Ours) & \textbf{74.20} & \textbf{69.47} & 63.83 & \textbf{29.00} \\
\bottomrule
\end{tabular}
\end{small}
\end{center}
\vskip -0.1in
\end{table*}

\subsection{Performance under Domain Shifts and Corruptions}
As shown in Table~\ref{tab:results}, under our rigorous deterministic protocol, our proposed frameworks decisively outperform every baseline across all environments, completing a perfect clean sweep. 
Under the standard clean environment, Static Merged achieves only 60.77\% accuracy due to severe parameter interference between the independently fine-tuned expert models. Active test-time adaptation (Standard TTA) mitigates this interference by learning optimal layer-wise weights, boosting performance to 74.00\%. Our proposed SATA-SBF and SATA-RGP further elevate this performance to \textbf{74.20\%}, indicating that enforcing convexity and sharpness constraints is beneficial even in-domain.

Under Gaussian Noise, the feature distributions undergo severe high-frequency corruptions. Standard TTA degrades to 67.10\%. Enforcing flatness via SAM TTA improves this to 69.27\% by preventing the model from converging into sharp minima that are highly sensitive to noise. Incorporating parameter-specific Fisher scaling (SBF TTA) further refines the performance to 69.30\%. Our unified SATA-RGP achieves the highest overall accuracy of \textbf{69.47\%}, representing an exceptional \textbf{+2.37\%} absolute increase over Standard TTA. This substantial improvement demonstrates the critical role of Relative Geometry Preservation, which ensures that the class prototype relationships are maintained even when the input features are corrupted.

Under Gaussian Blur, spatial patterns are severely smoothed, which heavily degrades early convolutional representations. Here, SATA-SBF leads with \textbf{63.87\%}, outperforming SAM TTA (63.73\%) and SBF TTA (63.67\%). 

Under Contrast reduction, the dynamic range of pixels is severely compressed, leading to extremely weak edge signals. This represents the most challenging domain shift, where the Static Merged model completely collapses to 16.30\% (only slightly above random guess). Standard TTA partially recovers to 28.17\%, but SAM TTA and SBF TTA fall back to 27.93\% and 27.97\%, respectively. This performance drop indicates that unconstrained optimization in a highly distorted feature space is extremely unstable and leads to representation collapse. In contrast, our proposed frameworks (SATA-SBF and SATA-RGP) boost accuracy to \textbf{29.00\%}. This represents a massive \textbf{+1.03\%} absolute improvement over the unconstrained SBF TTA, highlighting the vital regularizing effect of the Softmax constraint in Convex Tensor-wise Model Merging (C-TMM).

\begin{figure*}[t]
\centering
\includegraphics[width=0.85\textwidth]{tta_merging_results.png}
\caption{Visualization of multi-task test-time adaptation performance under various environmental domain shifts (Clean, Gaussian Noise, Gaussian Blur, and Contrast) across different merging methods.}
\label{fig:results_plot}
\end{figure*}

\subsection{Ablation Study}
To dissect the individual contributions of our core innovations, we conduct an exhaustive ablation study analyzing the effect of Convex TMM (C-TMM) via Softmax constraints versus unconstrained TMM across both the SATA-SBF and SATA-RGP configurations. The results are presented in Table~\ref{tab:ablation}.

\begin{table*}[t]
\caption{Ablation study analyzing the impact of Convex C-TMM (Softmax) vs. Unconstrained TMM on multi-task average accuracy (\%).}
\label{tab:ablation}
\vskip 0.15in
\begin{center}
\begin{small}
\begin{tabular}{lcccc}
\toprule
\textbf{Configuration} & \textbf{Clean} & \textbf{Gaussian Noise} & \textbf{Gaussian Blur} & \textbf{Contrast} \\
\midrule
SATA-SBF (Unconstrained TMM) & 74.07 & 69.27 & 63.77 & 27.97 \\
SATA-SBF (Convex C-TMM) & \textbf{74.20} & 69.37 & \textbf{63.87} & \textbf{29.00} \\
\midrule
SATA-RGP (Unconstrained TMM) & 74.07 & 69.40 & 63.70 & 27.97 \\
SATA-RGP (Convex C-TMM) & \textbf{74.20} & \textbf{69.47} & 63.83 & \textbf{29.00} \\
\bottomrule
\end{tabular}
\end{small}
\end{center}
\vskip -0.1in
\end{table*}

The ablation results reveal several profound insights:
\paragraph{Convex C-TMM consistently boosts performance:}
Comparing the unconstrained and convex configurations, we observe that C-TMM provides a consistent and substantial boost across all shifts. On the Clean stream, it improves both SATA-SBF and SATA-RGP from 74.07\% to \textbf{74.20\%}. Under Gaussian Noise, it increases SATA-SBF from 69.27\% to 69.37\% (+0.10\%) and SATA-RGP from 69.40\% to 69.47\% (+0.07\%). Under Gaussian Blur, it raises SATA-SBF from 63.77\% to 63.87\% (+0.10\%) and SATA-RGP from 63.70\% to 63.83\% (+0.13\%). Most decisively, under the severe Contrast shift, C-TMM boosts accuracy from 27.97\% to \textbf{29.00\%} (\textbf{+1.03\%} absolute increase). 

Mathematically, unconstrained TMM allows the layer-wise merging coefficients to take arbitrary real values, which can scale the task vectors excessively and lead to weight explosion or drift into invalid regions of the parameter space. By applying the Softmax constraint, C-TMM restricts the parameters to the convex hull of the expert trajectories, ensuring that the merged model's representations are always a valid interpolation of the well-trained experts. This mathematical boundary acts as a powerful regularizer, particularly when the incoming features are highly degraded.

\paragraph{RGP outperforms SPOR under Noise:}
When comparing SATA-RGP to SATA-SBF (which uses SPOR) under the unconstrained setting, SATA-RGP outperforms SATA-SBF by \textbf{+0.13\%} under Gaussian Noise (69.40\% vs 69.27\%) and maintains a similar margin of \textbf{+0.10\%} in the convex setting (69.47\% vs 69.37\%). 

This performance gap stems from the fundamental geometric difference between RGP and SPOR. SPOR enforces cross-correlation orthogonality, which implicitly assumes that the expert's classification head is perfectly orthogonal ($W^0 (W^0)^T = I$). In practical fine-tuned networks, however, classification boundaries are rarely perfectly orthogonal, as they naturally learn to reflect the semantic similarities between different classes (e.g., in KMNIST or FashionMNIST, certain classes share common visual patterns). Forcing the adapted classification head to satisfy a rigid orthogonality constraint distorts these learned semantic similarities, leading to decision-boundary degradation. In contrast, RGP aligns the adapted head's Gram matrix directly with the expert's Gram matrix, preserving the exact pairwise relative angles and cosine similarities between class prototypes. This allows the model to retain its optimal semantic classification boundaries under noise, demonstrating the superiority of relative geometry preservation over rigid orthogonalization.

\section{Discussion, Limitations, and Broader Impact}
\label{sec:discussion}

\subsection{Discussion and Scalability to Large Foundation Models}
A primary advantage of our proposed frameworks is their parameter-efficiency and computational viability. While standard TTA methods require full backward passes through the entire backbone network (which is extremely slow and memory-intensive for large foundation models like LLaMA \cite{touvron23llama} or Vision Transformers \cite{dosovitskiy21vit}), our approach limits active parameters to a small set of layer-wise merging coefficients (24 parameters) and the classification heads. Furthermore, the diagonal Fisher Information Matrix can be tracked online with negligible computational and memory overhead, as it only requires storing a running mean of the squared gradients of the active parameters. This makes SATA-SBF and SATA-RGP highly scalable and directly applicable to large-scale multi-task foundation models, enabling zero-shot, real-time adaptation of merged models on edge devices.

\subsection{Limitations}
Despite the strong empirical and theoretical foundation, our approach has several limitations:
\begin{enumerate}
    \item \textbf{Non-Stationary and Temporal Drift}: Our online TTA assumes that the incoming stream is composed of interleaved task batches. In extremely non-stationary environments where the task distribution changes abruptly or contains long periods of single-class streams, the running Fisher estimates can become biased, potentially leading to local overfitting. Introducing temporal smoothing or EMA buffers (such as DTS-TTA) could help mitigate this.
    \item \textbf{Dependence on Expert Weights}: The expert-guided self-labeling requires access to the original, frozen expert models during test-time to generate soft targets. While this is feasible in many settings, it increases the storage requirements at test-time since multiple expert encoders must be kept in memory, albeit frozen.
\end{enumerate}

\subsection{Broader Impact}
Our work has a positive broader impact on the environmental sustainability and democratization of AI. Retraining massive foundation models to integrate multiple specialized tasks is extremely energy-intensive and restricted to organizations with massive compute resources. Model merging provides a democratic, training-free alternative that can combine models at the weight level in seconds. By making model merging highly robust and adaptive to real-world corruptions and shifts, our frameworks pave the way for deploying unified, resource-efficient multi-task systems in challenging real-world environments, such as autonomous driving, robotics, and edge computing.

\section{Conclusion}
\label{sec:conclusion}
In this work, we introduced \textbf{SATA-SBF} and \textbf{SATA-RGP}, two robust online test-time adaptive model merging frameworks. We proposed Convex Tensor-wise Model Merging (C-TMM), which constrains layer-wise merging coefficients to the convex hull of experts via Softmax. We integrated Soft-Bounded Fisher-Guided SAM to scale sharpness updates inversely with online running Fisher sensitivities, shielding critical structures. Finally, we introduced Relative Geometry Preservation (RGP) to preserve the semantic similarity relationships of classification heads. Evaluated under a deterministic protocol, our frameworks decisively sweep all baselines across all environments, establishing a highly compatible, robust, and geometrically principled foundation for online adaptive model merging.

\clearpage

\nocite{*}
\bibliography{references}
\bibliographystyle{icml2026}

\end{document}
